% part: intuitionistic-logic
% chapter: propositions-as-types
% section: introduction

\documentclass[../../../include/open-logic-section]{subfiles}

\begin{document}

\olfileid{int}{pty}{int}

\olsection{પરિચય}

લૅમ્બડા-કલન ગણતરીનું એક સરળ નિદર્શ છે. તે
ચલોથી બનેલાં પદો પર કામ કરે છે. $N$ અને $M$
પદો હોય, તો $(NM)$ પણ એક પદ છે
(``$N$ને~$M$ પર લાગુ કર્યું''). $N$ પદ
હોય, તો $\lambd[x][N]$ પણ પદ છે. $N$માં
ચલ~$x$ આવતો હોય, તો $\lambd[x][N]$માં
$x$ની સંબંધિત આવૃત્તિઓ $\lambd[x]$ ક્રિયાકારકથી
બદ્ધ થાય છે અને હવે મુક્ત રહેતી નથી.
$(\lambd[x][N]M)$ સ્વરૂપનું પદ
$\Subst{N}{M}{x}$માં $\beta$-સંકોચાય છે
(એટલે~$N$માં $x$ની દરેક મુક્ત આવૃત્તિને
$M$થી બદલી તેનું પરિણામ મળે છે).
$N$ના ઉપપદને $\beta$-સંકોચવાથી $N'$
મળે, તો $N$ એ~$N'$માં $\beta$-રૂપાંતરિત
થાય છે; તેને $N \redone N'$ લખીએ.
$\lambd$-કલનમાં ગણતરી એ
$N \redone N_1 \redone N_2 \redone \dots$
અનુક્રમ છે. કોઈ ઉપપદ $\beta$-સંકોચી ન
શકાય એવા પદ~$N_k$ પર આ અનુક્રમ સમાપ્ત
થાય, તો તે પદ \emph{સામાન્ય} કહેવાય,
અથવા~$N$નું સામાન્ય સ્વરૂપ કહેવાય. $\lambd$-કલનની
ગણતરીને આવા અનુક્રમ તરીકે વિચારો. સામાન્ય
સ્વરૂપ મળે તો ગણતરી થોભે છે અને એ સ્વરૂપ
ગણતરીનું પરિણામ છે. ગણતરીનું આ સરળ
નિદર્શ ટ્યુરિંગ-પૂર્ણ છે: દરેક આંશિક પુનરાવર્તી
વિધેય $\lambd$-કલનના કોઈ પદથી ગણી શકાય છે.

હાસ્કેલ કરી અને વિલિયમ હાવર્ડે સ્વતંત્ર રીતે
સ્વાભાવિક નિગમન અને $\lambd$-કલન વચ્ચેનું
નજીકનું સામ્ય શોધ્યું. સહજ રીતે,
$\lambd[x][N]$ પદ એક વિધેય છે: $x$ તેનો
આર્ગ્યુમેન્ટ છે અને~$x$ આપવામાં આવે ત્યારે
મૂલ્ય કેવી રીતે ગણવું તે~$N$ કહે છે.
$(\lambd[x][N]M)$ પદ એટલે
$\lambd[x][N]$ વિધેયને આર્ગ્યુમેન્ટ~$M$
પર લાગુ કરવું; તેનું મૂલ્યાંકન કેવી રીતે કરવું
તે $\beta$-સંકોચન કહે છે: $N$માં $x$ને~$M$થી
બદલો (અને મળેલા પદનું મૂલ્યાંકન ચાલુ રાખો).
હવે આ વિધેયોનો પ્રદિશ અને મૂલ્ય-પ્રદેશ
નિશ્ચિત છે એમ વિચારો. $\lambd[x][N]$નો
પ્રદિશ~$A$ અને મૂલ્ય-પ્રદેશ~$B$ હોય, તો
ચલ~$x$ ફક્ત~$A$માંથી આર્ગ્યુમેન્ટ લે છે
અને~$N$ ફક્ત~$B$માંથી મૂલ્યો આપે છે.
એટલે $(\lambd[x][N]M)$ને અર્થ ત્યારે જ
છે જ્યારે $\lambd[x][N]$ એ $A \to B$
વિધેય હોય અને $M \in A$ હોય. આ માહિતી
$\lambd$-કલનમાં ઉમેરીએ, તો
\emph{પ્રકારિત} $\lambd$-કલન મળે છે.
પ્રકારિત $\lambd$-કલનમાં દરેક
$(\lambd[x][N]M)$ અભિવ્યક્તિ માન્ય નથી;
ફક્ત $\lambd[x][N]$ અને $M$ના
``પ્રકારો'' પરસ્પર બંધબેસે ત્યારે, એટલે
$\lambd[x][N]$નો પ્રકાર $A \lif B$
અને $M$નો પ્રકાર~$A$ હોય ત્યારે, માન્ય છે.
$x$નો પ્રકાર~$A$ અને $N$નો પ્રકાર~$B$
હોય તો જ $\lambd[x][N]$નો પ્રકાર
$A \to B$ છે. સામાન્ય રીતે દરેક
$(M_1M_2)$ અભિવ્યક્તિ માન્ય નથી.
$M_1$નો પ્રકાર $A \to B$ અને
$M_2$નો પ્રકાર~$A$ હોય એવા
$(M_1M_2)$ પદો જ માન્ય છે;
એમ હોય તો $(M_1M_2)$નો પ્રકાર~$B$ છે.

આ પ્રકાર-નિયમો સ્વાભાવિક નિગમનની
!!{derivation}sને સ્પષ્ટ રીતે અનુરૂપ છે:
પ્રકારોને !!{formula}s દર્શાવે છે અને
!!{derivation}s પોતે $\lambd$-પદોને
અનુરૂપ છે. દાખલા તરીકે,
$!A \lif !B$ની !!a{derivation}
$\lambd[\typeof{x}{!A}][N]$ પદને
અનુરૂપ છે; આ પદનો વિધેય-પ્રકાર
$!A \lif !B$ છે. સ્વાભાવિક નિગમનના
અનુમાન-નિયમો પ્રકારિત લૅમ્બડા-કલનના
પ્રકાર-નિયમોને અનુરૂપ છે. દાખલા તરીકે,
\begin{prooftree}
  \AxiomC{$\Discharge{!A}{x}$}
  \DeduceC{$!B$}
  \DischargeRule{\Intro\lif}{x}
  \UnaryInfC{$!A \lif !B$}
  \DisplayProof\bottomAlignProof
  \quad\text{ને અનુરૂપ}\quad
  \Axiom$x: !A \fCenter N: !B$
  \UnaryInf$ \fCenter \lambd[\typeof{x}{!A}][N] : !A \lif !B$
\end{prooftree}
અહીં જમણી બાજુનો નિયમ કહે છે કે $x$નો
પ્રકાર~$!A$ અને $N$નો પ્રકાર~$!B$
હોય, તો $\lambd[\typeof{x}{!A}][N]$નો
પ્રકાર~$!A \lif !B$ છે.

$\Elim\lif$ નિયમ સંયુક્ત પદો માટેના
પ્રકાર-નિયમને અનુરૂપ છે, એટલે
\[
  \AxiomC{$!A \lif !B$}
  \AxiomC{$!A$}
  \RightLabel{\Elim\lif}
  \BinaryInfC{$!B$}
  \DisplayProof
  \quad\text{ને અનુરૂપ}\quad
  \Axiom$\fCenter M_1: !A \lif !B$
  \Axiom$\fCenter M_2: !A$
  \BinaryInf$ \fCenter (M_1M_2) : !B$
  \DisplayProof
\]
\Intro\lif{} નિયમ પછી તરત \Elim\lif{}
નિયમ આવે, તો !!{derivation}ને સરળ બનાવી શકાય:
\begin{gather*}
  \AxiomC{$\Discharge{!A}{x}$}
  \DeduceC{$!B$}
  \DischargeRule{\Intro\lif}{x}
  \UnaryInfC{$!A \lif !B$}
  \AxiomC{}
  \DeduceC{$!A$}
  \RightLabel{\Elim\lif}
  \BinaryInfC{$!B$}
  \DisplayProof
  \quad\redone\quad
  \AxiomC{}
  \DeduceC{$!A$}
  \DeduceC{$!B$}
  \DisplayProof
\end{gather*}
આ લૅમ્બડા પદોના $\beta$-સંકોચનને અનુરૂપ છે:
\[
(\lambd[\typeof{x}{!A}][N]M) \quad\redone\quad
\Subst{N}{M}{x}.
\]

આવું જ અનુરૂપતા $\land$ના નિયમો અને
``ગુણાકાર'' પ્રકારો વચ્ચે, તથા
$\lor$ના નિયમો અને ``સરવાળા''
પ્રકારો વચ્ચે છે.

સરળ પ્રકારિત લૅમ્બડા-કલનના પદો અને
સ્વાભાવિક નિગમનની !!{derivation}s વચ્ચેની
આ અનુરૂપતાને ``કરી--હાવર્ડ'',
``પ્રોગ્રામ તરીકે સાબિતીઓ'' અથવા
``પ્રકારો તરીકે પ્રતિજ્ઞાપનો''ની
અનુરૂપતા કહે છે. !!{formula}s, એટલે
પ્રતિજ્ઞાપનો, પ્રકારોને અને સાબિતીઓ
પદોને અનુરૂપ છે; બાકીની અનુરૂપતાઓનો
સાર આ છે:
\begin{center}
  \begin{tabular}{c c c}
    તર્ક & પ્રોગ્રામ \\
    \hline
    પ્રતિજ્ઞાપન/!!{formula} & પ્રકાર \\
    સાબિતી & પદ \\
    ધારણા & ચલ \\
    મુક્ત કરેલી ધારણા & બદ્ધ ચલ \\
    મુક્ત ન કરેલી ધારણા & મુક્ત ચલ \\
    $!A \lif !B$ & વિધેય-પ્રકાર \\
    $!A \land !B$ & ગુણાકાર-પ્રકાર \\
    $!A \lor !B$ & સરવાળો-પ્રકાર \\
    $\lfalse$ & ખાલી પ્રકાર \\
    \hline
  \end{tabular}
\end{center}

કરી--હાવર્ડ અનુરૂપતા સ્વચાલિત સાબિતી-સહાયકો
અને પ્રોગ્રામના પ્રકાર-ચકાસકોનો એક આધારસ્તંભ
છે. કારણ કે, પ્રતિજ્ઞાપનની સાબિતી તપાસવી
(જેમ ઉપર કર્યું) એ પ્રોગ્રામ, એટલે પદ,
જાહેર કરેલા પ્રકારનો છે કે નહીં તે તપાસવા
બરાબર છે.
\end{document}
